\subsubsection{Ablation on Calibration Granularity}
\label{sec:ablation_granularity}

\begin{wraptable}[6]{r}{0.40\linewidth}
    \vspace{-1.6em}
    \centering
    \caption{Calibration granularity ablation.}
    \label{tab:ablation_granularity}
    \scriptsize
    \setlength{\tabcolsep}{5pt}
    \renewcommand{\arraystretch}{1.02}
    \begin{tabular}{lccc}
        \toprule
        \textbf{Metric} & \textbf{Tensor} & \textbf{Token} & \textbf{Channel} \\
        \midrule
        Coverage & 99.60\% & 99.74\% & 99.98\% \\
        Compression & $1.323\times$ & $1.322\times$ & $1.329\times$ \\
        Encode(GB/s) & 93.9 & 0.082 & 0.020 \\
        Decode(GB/s) & 222.2 & 0.173 & 0.060 \\
        \bottomrule
    \end{tabular}
    \vspace{-1.0em}
\end{wraptable}

Some prior compression designs use finer-grained codebooks to improve locality, compression ratio, or parallelism.
We evaluate this idea on per-tensor, per-token and per-channel codebooks.

Table~\ref{tab:ablation_granularity} shows that finer-grained calibration increases top-16 coverage, reaching 99.98\% with per-channel codebooks.
However, the projected compression ratio improves only marginally, from $1.323\times$ to $1.329\times$.
In contrast, throughput drops by several orders of magnitude: per-token and per-channel calibration both fall from hundreds of GB/s to sub-GB/s performance.
This is because fine-grained calibration requires many small codebooks, introducing extra storage, irregular lookup patterns, and more complex GPU execution.
Therefore, SplitZip uses a single global codebook, which preserves nearly the same compression ratio while maintaining a simple and high-throughput dense path.